\section{Institutional Background}
\label{sec:background}

\subsection{From Incurred Losses to Expected Credit Losses}
\label{sec:background:transition}

Under the historical incurred loss model (ILM), U.S.\ banks recognized
credit losses only when a loss became "probable and estimable,"
producing loan-loss provisions that lagged realized defaults
\citep{beatty2011delays} and reduced the informativeness of bank
disclosures \citep{bushman2012accounting, bushman2015delayed}.  Following
the 2007--2009 crisis, FASB issued ASU 2016-13, introducing the Current
Expected Credit Loss (CECL) model, codified in ASC~326
\citep{FASB2016}: banks estimate and reserve for \emph{lifetime}
expected credit losses at origination and at each reporting date.  The
required allowance is a direct function of loan-portfolio credit-risk
composition---longer maturities, weaker borrowers, and thinner collateral
require larger allowances than the ILM reserve.  Recent evidence finds
CECL adoption increased provision sensitivity to forward-looking credit
signals \citep{granja2025current, kim2026current}, reduced procyclicality
of bank lending \citep{chen2022does}, and improved decision-usefulness of
reserve disclosures \citep{gee2022decision}.

\subsection{The Day-One Adjustment and Treatment Variable}
\label{sec:background:dayone}

When a bank transitions from ILM to CECL, it must reclassify its entire
existing loan portfolio under the new standard on a single date.  The
difference between the CECL-required allowance and the legacy ILM allowance
(the Day-One adjustment) is recorded as a cumulative-effect charge to
retained earnings on the opening balance sheet of the adoption quarter.  For
bank $i$ with adoption quarter $t^{*}$:
\begin{equation}
  \Delta_i
  \;\equiv\;
  \mathrm{ALLL}^{\mathrm{CECL}}_{i,\,t^{*}}
  -
  \mathrm{ALLL}^{\mathrm{ILM}}_{i,\,t^{*}-1},
  \label{eq:dayone}
\end{equation}
where $\mathrm{ALLL}$ is the allowance for loan and lease losses on the Call
Report.  The Day-One adjustment varies widely: a large share of community
banks recorded a negligible adjustment, a subset recorded positive values
where CECL required a materially larger allowance, and a smaller number
recorded negative values.

Where positive, the adjustment reduces reported book equity and, absent
regulatory accommodation, risk-weighted capital ratios by an equivalent
amount.  It is a pure accounting event: cash flows and borrower repayment
terms are unchanged.  What the adjustment reveals, in a single mandatory
disclosure, is the gap between book equity as previously reported and book
equity adjusted for the bank's forward-looking credit-loss exposure.

Our primary treatment scales the Day-One charge by beginning-of-period
equity:
\begin{equation}
  \texttt{CECL Adj}_{i}
  \;=\;
  \frac{\Delta_i}{\texttt{Equity}_{i,\,t-1^{*}}}
  \times 100,
  \label{eq:treatment}
\end{equation}
expressed in percentage points.  For binary specifications we construct
\begin{equation}
  \texttt{High CECL}_{i}
  \;=\;
  \mathbf{1}\!\left\{
    \texttt{CECL Adj}_{i} \geq \tau
  \right\},
  \label{eq:binary}
\end{equation}
with $\tau = 2.5$ percentage points (approximately the 95th percentile).
Results are robust to using the continuous measure, to alternative
percentile cutoffs, and to scaling by assets or loans instead of equity
(Section~\ref{sec:robustness}).  Equity scaling is our baseline because it
measures the fraction of the loss-absorbing buffer that the new standard
revealed to be already committed against expected credit losses.

\subsection{Adoption Timeline}
\label{sec:background:timeline}

Large SEC-filing registrants adopted CECL in 2020Q1; non-public
entities---the community banks in our main sample---adopted in 2023Q1.
Call Reports are posted publicly by the FFIEC and FDIC roughly thirty
days after quarter-end, so 2023Q1 Day-One adjustments became publicly
observable in late April and May 2023---after the March 2023 failures of
Silicon Valley Bank and Signature Bank.  Depositors reacting to the
disclosure did so against an information backdrop in which the March run
had already occurred, a sequencing we exploit in
Section~\ref{sec:identification}.  The 2020 large-bank cohort coincides
with the COVID-19 onset, which precludes using it to study deposit
responses.
